%=============================================================================!
% 2.S  WHAT WE NOW HAVE                                                       !
%=============================================================================!
\unnumbered{What we now have}

A force is a push or a pull with a size and a direction, and the forces
on a body add as vectors to a net force. Accelerations are the same in
all frames of reference moving at constant velocity relative to one
another, and in the inertial frames among them, to which a frame fixed
to the ground is a close approximation, three laws hold: a body with no net force on it keeps
its velocity; the net force equals the mass times the acceleration; and
the forces between two bodies are equal and opposite. Under its weight
alone every body falls with the same acceleration $g$. The second law
turns every motion into a differential equation, whose solution is fixed
by the state, the position and velocity at one instant. With a drag
proportional to its velocity a body approaches its terminal velocity
$v_\infty = g/\gamma$ exponentially, the gap shrinking to
\SI{37}{\percent} every $1/\gamma$; on a spring it oscillates as a
sinusoid with $\omega = \sqrt{k/m}$, whatever its amplitude, as long as
Hooke's law holds. For several
bodies the internal forces cancel in pairs, so the total momentum changes
only through external forces and the centre of mass moves as a single
body. For atoms the same laws give $3N$ coupled equations, which in eV,
\AA, fs and amu need the single factor \num{9.648533e-3}, the
acceleration in \si{\angstrom\per\femto\second\squared} that a force of
\SI{1}{\electronvolt\per\angstrom} gives a mass of \SI{1}{\amu}.

What remains is to understand forces through energy. Chapter~3 shows
which forces conserve the total energy of a motion, and why that matters
for every simulation.
